\documentclass[10pt]{extarticle}
\usepackage[a4paper,margin=1.5cm]{geometry}
\usepackage[T1]{fontenc}
\usepackage{amsmath,amssymb}
\usepackage{graphicx}
\usepackage[export]{adjustbox}
\usepackage{multicol}
\usepackage{enumitem}
\usepackage{booktabs,array,tabularx}
\setlength{\columnsep}{0.7cm}
\setlength{\parindent}{0pt}
\setlength{\parskip}{3pt}
\pagestyle{empty}
\makeatletter
% figure registry (figdefs.tex is generated by the extract step)
\newcommand{\qfigdef}[4]{\@namedef{qf@file@#1}{#2}\@namedef{qf@w@#1}{#3}\@namedef{qf@wide@#1}{#4}}
\newcommand{\qf@check}[1]{\@ifundefined{qf@file@#1}{\PackageError{score2}{Unknown figure id #1}{}}{}}
\newcommand{\qf@img}[4][]{\includegraphics[width=#2\dimexpr\@nameuse{qf@w@#4}pt\relax,max width=#3,max height=0.45\textheight,#1]{\@nameuse{qf@file@#4}}}
% Full-width questions: finalize wraps every question that contains a wide figure in \qfullbegin...\qfullend,
% so the WHOLE question (stem, figures, options) is set across the page in reading order, then the two columns
% resume. The question's head - its number, stem and first figure (\qfullsplit marks the end of the head) -
% is kept on one page, so a number is never left at a page foot with its figure overleaf; the rest may break
% across pages normally, so no big gaps are left. Without a \qfullsplit, the whole question is kept together
% if it is under 45% of a page. Never write these yourself.
\newif\ifqf@full
\newif\ifqf@boxed
\newsavebox{\qf@box}
\newcommand{\qf@open}{\global\qf@fulltrue\par\global\qf@boxedtrue
  \setbox\qf@box\vbox\bgroup\hsize\textwidth\linewidth\textwidth}
\newcommand{\qf@emit}[1]{\par\egroup\global\qf@boxedfalse
  \ifdim\dimexpr\ht\qf@box+\dp\qf@box\relax<#1\textheight\noindent\usebox{\qf@box}\par\else\unvbox\qf@box\fi}
\newcommand{\qfullsplit}{\ifqf@boxed\qf@emit{0.85}\fi}
\newcommand{\qf@close}{\ifqf@boxed\qf@emit{0.45}\fi\global\qf@fullfalse}
\newcommand{\qfullbegin}{\end{multicols}\qf@open}
\newcommand{\qfullend}{\qf@close\begin{multicols}{2}}
\newcommand{\qfullstart}{\qf@open}
\newcommand{\qfulldone}{\qf@close}
% \qfig[scale]{id}: block figure, centred, directly where it is written (size set by the tool).
% Fallback if a wide figure is not inside a full-width question: break out of the columns in place - no floats.
\newcommand{\qfig}[2][1]{\qf@check{#2}%
  \ifqf@full\qf@blk{#1}{#2}\else
  \ifnum\@nameuse{qf@wide@#2}=1 \expandafter\@firstoftwo\else\expandafter\@secondoftwo\fi
  {\end{multicols}\begin{center}\qf@img{#1}{\textwidth}{#2}\end{center}\begin{multicols}{2}}%
  {\qf@blk{#1}{#2}}\fi}
\newcommand{\qf@blk}[2]{\par\smallskip{\centering\qf@img{#1}{\linewidth}{#2}\par}\smallskip}
% \qfigi[scale]{id}: inline figure (inside option lists / tables / minipages); never breaks columns.
\newcommand{\qfigi}[2][1]{\qf@check{#2}\qf@img[valign=c]{#1}{\linewidth}{#2}}
\makeatother
% structure macros - lint reads these to check completeness and figure placement
\newcommand{\papersec}[1]{\par\bigskip{\centering\Large\bfseries PAPER-#1\par}\medskip}
\newcommand{\partsec}[2]{\par\medskip\noindent{\large\bfseries PART-#1: #2}\par\smallskip}
\newcommand{\sect}[1]{\par\smallskip\noindent\textbf{#1}\par}
\newcommand{\qn}[1]{\par\smallskip\noindent\textbf{#1.}~}
\newcommand{\paraq}[2]{\par\smallskip\noindent\textbf{Paragraph for Questions #1 and #2}\par}
\newenvironment{qoptions}{\begin{itemize}[leftmargin=1.8em,itemsep=1pt,topsep=2pt,label={}]}{\end{itemize}}
\newcommand{\opt}[2]{\item[(#1)] #2}
% figure registry (generated by score2_tools.py finalize)
\qfigdef{p1_phy_s1_q01_1}{TNAS - Test 01 (05 Feb 2023)-img/p1_phy_s1_q01_1.png}{108.3}{0}
\qfigdef{p1_phy_s1_q05_1}{TNAS - Test 01 (05 Feb 2023)-img/p1_phy_s1_q05_1.png}{108.3}{0}
\qfigdef{p1_phy_s1_q11_1}{TNAS - Test 01 (05 Feb 2023)-img/p1_phy_s1_q11_1.png}{108.2}{0}
\qfigdef{p1_phy_s1_q13_1}{TNAS - Test 01 (05 Feb 2023)-img/p1_phy_s1_q13_1.png}{108.2}{0}
\qfigdef{p1_phy_s1_q14_1}{TNAS - Test 01 (05 Feb 2023)-img/p1_phy_s1_q14_1.png}{55.6}{0}
\qfigdef{p1_phy_s1_q15_1}{TNAS - Test 01 (05 Feb 2023)-img/p1_phy_s1_q15_1.png}{87.2}{0}
\qfigdef{p1_phy_s1_q18_1}{TNAS - Test 01 (05 Feb 2023)-img/p1_phy_s1_q18_1.png}{64.7}{0}
\qfigdef{p1_phy_s1_q20_1}{TNAS - Test 01 (05 Feb 2023)-img/p1_phy_s1_q20_1.png}{79.4}{0}
\qfigdef{p1_phy_s2_q04_1}{TNAS - Test 01 (05 Feb 2023)-img/p1_phy_s2_q04_1.png}{154.5}{0}
\qfigdef{p1_phy_s2_q05_1}{TNAS - Test 01 (05 Feb 2023)-img/p1_phy_s2_q05_1.png}{154.2}{0}
\qfigdef{p1_phy_s2_q07_1}{TNAS - Test 01 (05 Feb 2023)-img/p1_phy_s2_q07_1.png}{124.1}{0}
\qfigdef{p1_phy_s2_q08_1}{TNAS - Test 01 (05 Feb 2023)-img/p1_phy_s2_q08_1.png}{63.5}{0}
\qfigdef{p1_phy_s2_q09_1}{TNAS - Test 01 (05 Feb 2023)-img/p1_phy_s2_q09_1.png}{108.2}{0}
\qfigdef{p1_phy_s2_q10_1}{TNAS - Test 01 (05 Feb 2023)-img/p1_phy_s2_q10_1.png}{108.3}{0}
\qfigdef{p1_chem_s1_q05_A}{TNAS - Test 01 (05 Feb 2023)-img/p1_chem_s1_q05_A.png}{30.8}{0}
\qfigdef{p1_chem_s1_q05_B}{TNAS - Test 01 (05 Feb 2023)-img/p1_chem_s1_q05_B.png}{43.1}{0}
\qfigdef{p1_chem_s1_q12_A}{TNAS - Test 01 (05 Feb 2023)-img/p1_chem_s1_q12_A.png}{71.5}{0}
\qfigdef{p1_chem_s1_q12_B}{TNAS - Test 01 (05 Feb 2023)-img/p1_chem_s1_q12_B.png}{71.0}{0}
\qfigdef{p1_chem_s1_q12_C}{TNAS - Test 01 (05 Feb 2023)-img/p1_chem_s1_q12_C.png}{71.3}{0}
\qfigdef{p1_chem_s1_q12_D}{TNAS - Test 01 (05 Feb 2023)-img/p1_chem_s1_q12_D.png}{71.1}{0}
\qfigdef{p1_chem_s1_q14_A}{TNAS - Test 01 (05 Feb 2023)-img/p1_chem_s1_q14_A.png}{136.4}{0}
\qfigdef{p1_chem_s1_q14_2}{TNAS - Test 01 (05 Feb 2023)-img/p1_chem_s1_q14_2.png}{35.4}{0}
\qfigdef{p1_chem_s1_q14_3}{TNAS - Test 01 (05 Feb 2023)-img/p1_chem_s1_q14_3.png}{21.3}{0}
\qfigdef{p1_chem_s1_q14_4}{TNAS - Test 01 (05 Feb 2023)-img/p1_chem_s1_q14_4.png}{29.2}{0}
\qfigdef{p1_chem_s1_q14_D}{TNAS - Test 01 (05 Feb 2023)-img/p1_chem_s1_q14_D.png}{137.4}{0}
\qfigdef{p1_chem_s1_q15_1}{TNAS - Test 01 (05 Feb 2023)-img/p1_chem_s1_q15_1.png}{69.3}{0}
\qfigdef{p1_chem_s1_q15_2}{TNAS - Test 01 (05 Feb 2023)-img/p1_chem_s1_q15_2.png}{85.2}{0}
\qfigdef{p1_chem_s1_q15_3}{TNAS - Test 01 (05 Feb 2023)-img/p1_chem_s1_q15_3.png}{66.6}{0}
\qfigdef{p1_chem_s1_q16_1}{TNAS - Test 01 (05 Feb 2023)-img/p1_chem_s1_q16_1.png}{132.0}{0}
\qfigdef{p1_chem_s1_q16_A}{TNAS - Test 01 (05 Feb 2023)-img/p1_chem_s1_q16_A.png}{37.2}{0}
\qfigdef{p1_chem_s1_q16_B}{TNAS - Test 01 (05 Feb 2023)-img/p1_chem_s1_q16_B.png}{37.2}{0}
\qfigdef{p1_chem_s1_q16_C}{TNAS - Test 01 (05 Feb 2023)-img/p1_chem_s1_q16_C.png}{37.0}{0}
\qfigdef{p1_chem_s1_q16_D}{TNAS - Test 01 (05 Feb 2023)-img/p1_chem_s1_q16_D.png}{37.0}{0}
\qfigdef{p1_chem_s1_q17_1}{TNAS - Test 01 (05 Feb 2023)-img/p1_chem_s1_q17_1.png}{76.2}{0}
\qfigdef{p1_chem_s1_q17_A}{TNAS - Test 01 (05 Feb 2023)-img/p1_chem_s1_q17_A.png}{24.0}{0}
\qfigdef{p1_chem_s1_q17_B}{TNAS - Test 01 (05 Feb 2023)-img/p1_chem_s1_q17_B.png}{21.5}{0}
\qfigdef{p1_chem_s1_q17_C}{TNAS - Test 01 (05 Feb 2023)-img/p1_chem_s1_q17_C.png}{21.3}{0}
\qfigdef{p1_chem_s1_q18_1}{TNAS - Test 01 (05 Feb 2023)-img/p1_chem_s1_q18_1.png}{43.0}{0}
\qfigdef{p1_chem_s1_q18_2}{TNAS - Test 01 (05 Feb 2023)-img/p1_chem_s1_q18_2.png}{53.2}{0}
\qfigdef{p1_chem_s1_q18_B}{TNAS - Test 01 (05 Feb 2023)-img/p1_chem_s1_q18_B.png}{122.0}{0}
\qfigdef{p1_chem_s1_q18_C}{TNAS - Test 01 (05 Feb 2023)-img/p1_chem_s1_q18_C.png}{47.5}{0}
\qfigdef{p1_chem_s1_q18_5}{TNAS - Test 01 (05 Feb 2023)-img/p1_chem_s1_q18_5.png}{55.6}{0}
\qfigdef{p1_chem_s1_q18_D}{TNAS - Test 01 (05 Feb 2023)-img/p1_chem_s1_q18_D.png}{76.2}{0}
\qfigdef{p1_chem_s1_q18_7}{TNAS - Test 01 (05 Feb 2023)-img/p1_chem_s1_q18_7.png}{31.8}{0}
\qfigdef{p1_chem_s1_q19_A}{TNAS - Test 01 (05 Feb 2023)-img/p1_chem_s1_q19_A.png}{125.8}{0}
\qfigdef{p1_chem_s1_q19_B}{TNAS - Test 01 (05 Feb 2023)-img/p1_chem_s1_q19_B.png}{135.2}{0}
\qfigdef{p1_chem_s1_q19_C}{TNAS - Test 01 (05 Feb 2023)-img/p1_chem_s1_q19_C.png}{128.3}{0}
\qfigdef{p1_chem_s1_q19_D}{TNAS - Test 01 (05 Feb 2023)-img/p1_chem_s1_q19_D.png}{133.7}{0}
\qfigdef{p1_chem_s1_q20_1}{TNAS - Test 01 (05 Feb 2023)-img/p1_chem_s1_q20_1.png}{142.8}{0}
\qfigdef{p1_chem_s1_q20_2}{TNAS - Test 01 (05 Feb 2023)-img/p1_chem_s1_q20_2.png}{44.9}{0}
\qfigdef{p1_chem_s2_q08_1}{TNAS - Test 01 (05 Feb 2023)-img/p1_chem_s2_q08_1.png}{128.6}{0}
\qfigdef{p1_chem_s2_q08_2}{TNAS - Test 01 (05 Feb 2023)-img/p1_chem_s2_q08_2.png}{140.4}{0}
\qfigdef{p1_chem_s2_q08_3}{TNAS - Test 01 (05 Feb 2023)-img/p1_chem_s2_q08_3.png}{131.5}{0}
\qfigdef{p1_chem_s2_q08_4}{TNAS - Test 01 (05 Feb 2023)-img/p1_chem_s2_q08_4.png}{74.5}{0}
\qfigdef{p1_chem_s2_q09_1}{TNAS - Test 01 (05 Feb 2023)-img/p1_chem_s2_q09_1.png}{50.6}{0}
\qfigdef{p1_chem_s2_q09_2}{TNAS - Test 01 (05 Feb 2023)-img/p1_chem_s2_q09_2.png}{50.6}{0}
\qfigdef{p1_chem_s2_q09_3}{TNAS - Test 01 (05 Feb 2023)-img/p1_chem_s2_q09_3.png}{66.3}{0}
\qfigdef{p1_chem_s2_q09_4}{TNAS - Test 01 (05 Feb 2023)-img/p1_chem_s2_q09_4.png}{43.1}{0}
\qfigdef{p1_chem_s2_q09_5}{TNAS - Test 01 (05 Feb 2023)-img/p1_chem_s2_q09_5.png}{50.6}{0}
\qfigdef{p1_chem_s2_q09_6}{TNAS - Test 01 (05 Feb 2023)-img/p1_chem_s2_q09_6.png}{60.8}{0}
\qfigdef{p1_chem_s2_q10_1}{TNAS - Test 01 (05 Feb 2023)-img/p1_chem_s2_q10_1.png}{50.2}{0}
\qfigdef{p2_phy_s1_q02_1}{TNAS - Test 01 (05 Feb 2023)-img/p2_phy_s1_q02_1.png}{157.9}{0}
\qfigdef{p2_phy_s1_q04_1}{TNAS - Test 01 (05 Feb 2023)-img/p2_phy_s1_q04_1.png}{132.2}{0}
\qfigdef{p2_phy_s1_q05_1}{TNAS - Test 01 (05 Feb 2023)-img/p2_phy_s1_q05_1.png}{68.9}{0}
\qfigdef{p2_phy_s1_q06_1}{TNAS - Test 01 (05 Feb 2023)-img/p2_phy_s1_q06_1.png}{124.1}{0}
\qfigdef{p2_phy_s1_q08_1}{TNAS - Test 01 (05 Feb 2023)-img/p2_phy_s1_q08_1.png}{82.0}{0}
\qfigdef{p2_phy_s2_q09_1}{TNAS - Test 01 (05 Feb 2023)-img/p2_phy_s2_q09_1.png}{150.8}{0}
\qfigdef{p2_phy_s2_q10_1}{TNAS - Test 01 (05 Feb 2023)-img/p2_phy_s2_q10_1.png}{150.8}{0}
\qfigdef{p2_phy_s2_q11_1}{TNAS - Test 01 (05 Feb 2023)-img/p2_phy_s2_q11_1.png}{124.2}{0}
\qfigdef{p2_phy_s2_q12_1}{TNAS - Test 01 (05 Feb 2023)-img/p2_phy_s2_q12_1.png}{134.7}{0}
\qfigdef{p2_phy_s3_q01_1}{TNAS - Test 01 (05 Feb 2023)-img/p2_phy_s3_q01_1.png}{84.6}{0}
\qfigdef{p2_phy_s3_q02_1}{TNAS - Test 01 (05 Feb 2023)-img/p2_phy_s3_q02_1.png}{92.0}{0}
\qfigdef{p2_phy_s3_q03_1}{TNAS - Test 01 (05 Feb 2023)-img/p2_phy_s3_q03_1.png}{174.2}{0}
\qfigdef{p2_phy_s3_q04_1}{TNAS - Test 01 (05 Feb 2023)-img/p2_phy_s3_q04_1.png}{100.1}{0}
\qfigdef{p2_chem_s1_q03_A}{TNAS - Test 01 (05 Feb 2023)-img/p2_chem_s1_q03_A.png}{97.7}{0}
\qfigdef{p2_chem_s1_q03_B}{TNAS - Test 01 (05 Feb 2023)-img/p2_chem_s1_q03_B.png}{124.9}{0}
\qfigdef{p2_chem_s1_q03_C}{TNAS - Test 01 (05 Feb 2023)-img/p2_chem_s1_q03_C.png}{84.6}{0}
\qfigdef{p2_chem_s1_q03_D}{TNAS - Test 01 (05 Feb 2023)-img/p2_chem_s1_q03_D.png}{82.1}{0}
\qfigdef{p2_chem_s1_q04_A}{TNAS - Test 01 (05 Feb 2023)-img/p2_chem_s1_q04_A.png}{65.6}{0}
\qfigdef{p2_chem_s1_q04_B}{TNAS - Test 01 (05 Feb 2023)-img/p2_chem_s1_q04_B.png}{70.3}{0}
\qfigdef{p2_chem_s1_q04_C}{TNAS - Test 01 (05 Feb 2023)-img/p2_chem_s1_q04_C.png}{70.8}{0}
\qfigdef{p2_chem_s1_q04_D}{TNAS - Test 01 (05 Feb 2023)-img/p2_chem_s1_q04_D.png}{70.3}{0}
\qfigdef{p2_chem_s1_q05_A}{TNAS - Test 01 (05 Feb 2023)-img/p2_chem_s1_q05_A.png}{100.4}{0}
\qfigdef{p2_chem_s1_q05_B}{TNAS - Test 01 (05 Feb 2023)-img/p2_chem_s1_q05_B.png}{95.1}{0}
\qfigdef{p2_chem_s1_q05_C}{TNAS - Test 01 (05 Feb 2023)-img/p2_chem_s1_q05_C.png}{97.7}{0}
\qfigdef{p2_chem_s1_q05_D}{TNAS - Test 01 (05 Feb 2023)-img/p2_chem_s1_q05_D.png}{97.9}{0}
\qfigdef{p2_chem_s1_q06_B}{TNAS - Test 01 (05 Feb 2023)-img/p2_chem_s1_q06_B.png}{107.7}{0}
\qfigdef{p2_chem_s1_q06_C}{TNAS - Test 01 (05 Feb 2023)-img/p2_chem_s1_q06_C.png}{201.5}{0}
\qfigdef{p2_chem_s1_q06_D}{TNAS - Test 01 (05 Feb 2023)-img/p2_chem_s1_q06_D.png}{173.1}{0}
\qfigdef{p2_chem_s3_q02_1}{TNAS - Test 01 (05 Feb 2023)-img/p2_chem_s3_q02_1.png}{128.1}{0}
\raggedcolumns
\makeatletter\@beginparpenalty=10000\makeatother
\newcolumntype{L}[1]{>{\raggedright\arraybackslash}p{#1}}
\begin{document}
\begin{center}{\Large\bfseries TNAS --- Test 01 (05 Feb 2023)}\end{center}
\begin{multicols}{2}
\raggedcolumns
\papersec{1}
\partsec{1}{PHYSICS}
\sect{SECTION-I}
\qn{1} By what angle should $M_2$ be rotated, so that the light ray after reflection from both the mirrors become horizontal :-
\qfig{p1_phy_s1_q01_1}
\begin{qoptions}
\opt{A}{$5^\circ$ clockwise}
\opt{B}{$95^\circ$ clockwise}
\opt{C}{$5^\circ$ anticlockwise}
\opt{D}{$15^\circ$ anticlockwise}
\end{qoptions}
\qn{2} A ray of light moving along vector $\hat i-2\hat j$ in denser medium undergoes refraction at an interface of two media, which is $x$-$z$ plane. The refractive index for $y>0$ is $2$ while for $y<0$ it is $\dfrac{\sqrt5}{2}$. The unit vector along the refracted ray is :-
\begin{qoptions}
\opt{A}{$\dfrac{-3\hat i-5\hat j}{\sqrt{34}}$}
\opt{B}{$\dfrac{-4\hat i+5\hat j}{5}$}
\opt{C}{$\dfrac{-3\hat i-4\hat j}{5}$}
\opt{D}{$\dfrac{4\hat i-3\hat j}{5}$}
\end{qoptions}
\qn{3} A bird is flying over a swimming pool at a height of $2$ m from the water surface. If the bottom is perfectly plane reflecting surface and depth of swimming pool is $1$m, then the distance of final image of bird from the bird itself is $\mu_w=4/3$.
\begin{qoptions}
\opt{A}{$\dfrac{11}{3}$m}
\opt{B}{$\dfrac{23}{3}$m}
\opt{C}{$\dfrac{11}{4}$m}
\opt{D}{$\dfrac{11}{2}$m}
\end{qoptions}
\qn{4} A plane mirror is moving with velocity $4\hat i+4\hat j+8\hat k$. A point object in front of the mirror moves with a velocity $3\hat i+4\hat j+5\hat k$. Here, $\hat k$ is along the normal to the plane mirror and facing towards the object. The velocity of the image is :-
\begin{qoptions}
\opt{A}{$-3\hat i-4\hat j+5\hat k$}
\opt{B}{$3\hat i+4\hat j+11\hat k$}
\opt{C}{$-4\hat i+5\hat j+11\hat k$}
\opt{D}{$7\hat i+9\hat j+3\hat k$}
\end{qoptions}
\qn{5} A leakage begins in water tank at position P as shown in figure. The initial gauge pressure at P was $5\times10^5$ N/m$^2$. If the density of water is $1000$ kg/m$^3$ the initial velocity with which water gushes out is approximately ($P_{atm}=1\times10^5$)
\qfig{p1_phy_s1_q05_1}
\begin{qoptions}
\opt{A}{$3.2$ ms$^{-1}$}
\opt{B}{$35$ ms$^{-1}$}
\opt{C}{$32$ m/s}
\opt{D}{$3.5$ m/s$^{-1}$}
\end{qoptions}
\qn{6} An open container of height $H$ and radius $R$ is completely filled with water. A small orifice of radius $r$ is made near the bottom of container. Time elapsed in emptying half of the container, is
\begin{qoptions}
\opt{A}{$\dfrac{R}{r}\sqrt{\dfrac{2H}{g}}$}
\opt{B}{$\dfrac{R}{r}\sqrt{\dfrac{2H}{g}}\,(\sqrt2-1)$}
\opt{C}{$\dfrac{R^2}{r^2}\sqrt{\dfrac{2H}{g}}\left(1-\dfrac{1}{\sqrt2}\right)$}
\opt{D}{$\dfrac{R^2}{r^2}\sqrt{\dfrac Hg}$}
\end{qoptions}
\qn{7} A frustum shaped rod has cross sectional area that varies from its one end with distance $x$ as $A=A_0(1+\alpha x)$, where $A_0$ and $\alpha$ are constants. The rod has length $\ell$ and is maintained between the temperatures $T_1$ and $T_2$ ($T_1<T_2$) with $T_1$ being the temperature at cross section with smaller radius. Which of the following represents the temperature at a cross section at distance $x$ from the end at temperature $T_1$ ?
\begin{qoptions}
\opt{A}{$T_1+(T_2-T_1)\dfrac{\ell n\,(1+\alpha x)}{\ell n\,(1+\alpha\ell)}$}
\opt{B}{$T_1-T_2\dfrac{\ell n\,(1+\alpha x)}{\ell n\,(1+\alpha\ell)}$}
\opt{C}{$T_1-T_2\dfrac{\ell n\,(\alpha x)}{\ell n\,(\alpha\ell)}$}
\opt{D}{$T_2-T_1\dfrac{\ell n\,(1+\alpha x)}{\ell n\,(1+\alpha\ell)}$}
\end{qoptions}
\qn{8} Which among the following represents the slope of adiabatic process of an ideal gas in $\ell n T$ Vs $\ell n P$ graph.
\begin{qoptions}
\opt{A}{$\dfrac{C_P}{C_V}$}
\opt{B}{$1-\dfrac{C_P}{C_V}$}
\opt{C}{$1-\dfrac{C_V}{C_P}$}
\opt{D}{$-\dfrac{C_V}{C_P}$}
\end{qoptions}
\qn{9} A particle of mass $m$ confined to move along $x$-axis has potential energy given by $U(x)=\dfrac{a}{b^4}\left(x^4-2x^2b^2\right)$, where $a$ \& $b$ are positive constants and $x$ is position of particle. Frequency of small oscillation near the equilibrium point will be :-
\begin{qoptions}
\opt{A}{$\dfrac{1}{2\pi}\sqrt{\dfrac{8a}{mb^2}}$}
\opt{B}{$\dfrac{1}{2\pi}\sqrt{\dfrac{4a}{mb^2}}$}
\opt{C}{$2\pi\sqrt{\dfrac{4a}{mb^2}}$}
\opt{D}{$\pi\sqrt{\dfrac{2a}{mb^2}}$}
\end{qoptions}
\qn{10} A point mass at rest is subjected to three SHM displacements simultaneously in same direction. If equations of three SHMs are given by $x_1=A\sin\omega t$, $x_2=A\sin\left(\omega t+\dfrac{2\pi}{3}\right)$ and $x_3=B\sin(\omega t+\phi)$, then for mass to remain at rest.
\begin{qoptions}
\opt{A}{$B=A$, $\phi=\dfrac{4\pi}{3}$}
\opt{B}{$B=2A$, $\phi=\dfrac{\pi}{3}$}
\opt{C}{$B=A$, $\phi=\dfrac{\pi}{3}$}
\opt{D}{$B=\sqrt3A$, $\phi=\dfrac{5\pi}{6}$}
\end{qoptions}
\qn{11} The given plot shows the variation of $U$, the potential energy of interaction between two particles with the distance separating them $r$.
\qfig{p1_phy_s1_q11_1}

1. B and D are equilibrium points

2. C is a point of stable equilibrium

3. The force of interaction between the two particles is attractive between points C and D and repulsive between D and E.

4. The force of interaction between particles is repulsive between points E and F.

Which of the above statements are correct ?
\begin{qoptions}
\opt{A}{1 and 2}
\opt{B}{1 and 4}
\opt{C}{2 and 4}
\opt{D}{2 and 3}
\end{qoptions}
\qn{12} A rod of length $1000$ mm and coefficient of linear expansion $\alpha=10^{-4}$ per degree celsius is placed in horizontal smooth surface symmetrically between fixed walls separated by $1001$ mm. The Young's modulus of rod is $10^{11}$ N/m$^2$. If the temperature is increased by $20^\circ$C, then the stress developed in the rod is (in N/m$^2$)
\begin{qoptions}
\opt{A}{$10^5$}
\opt{B}{$10^8$}
\opt{C}{$10^7$}
\opt{D}{$10^6$}
\end{qoptions}
\qn{13} An ideal liquid is flowing through a pipe kept on horizontal floor as shown in figure. The pressure at two points are indicated, then
\qfig{p1_phy_s1_q13_1}
\begin{qoptions}
\opt{A}{$P_1>P_2$}
\opt{B}{$P_1<P_2$}
\opt{C}{$P_1=P_2$}
\opt{D}{$P_1$ and $P_2$ cannot be compared}
\end{qoptions}
\qn{14} A cylindrical open tank of radius $R$ and height $h$ is filled completely with water. It is now rotated about its vertical axis. The minimum angular speed such that the small part of bottom surface get exposed to atmosphere is
\qfig{p1_phy_s1_q14_1}
\begin{qoptions}
\opt{A}{$\dfrac{\sqrt{gh}}{R}$}
\opt{B}{$\dfrac{\sqrt{gh}}{2R}$}
\opt{C}{$\dfrac{2\sqrt{gh}}{R}$}
\opt{D}{$\dfrac{\sqrt{2gh}}{R}$}
\end{qoptions}
\qn{15} A fixed cylindrical tank having cross sectional area $A$ is filled with two liquids of density $\rho$, $2\rho$ to a height $h$, $2h$ respectively. A small hole having area $a$ is made in vertical wall so as to achieve maximum range of effluxed liquid on ground. If ($\rho=1000$ kg/m$^3$, $h=0.5$ m, $A=1$m$^2$, $a=1$ cm$^2$). Maximum range of liquid efflux on ground is :-
\qfig{p1_phy_s1_q15_1}
\begin{qoptions}
\opt{A}{$0.75$ m}
\opt{B}{$1$m}
\opt{C}{$1.25$ m}
\opt{D}{$\sqrt2$ m}
\end{qoptions}
\qn{16} Temperature at which rms speed of $\mathrm{SO_2}$ molecules is same as that of $\mathrm{O_2}$ molecules at $27^\circ$C, is (Given molecular weight of $\mathrm{O_2}$ and $\mathrm{SO_2}$ are $32$ and $64$ respectively)
\begin{qoptions}
\opt{A}{$327^\circ$C}
\opt{B}{$600^\circ$C}
\opt{C}{$227^\circ$C}
\opt{D}{$527^\circ$C}
\end{qoptions}
\qn{17} A particle executes simple harmonic motion with an amplitude of $5$cm. When the particle is at $4$ cm from the mean position, the magnitude of its velocity in SI units is equal to that of its acceleration. Then, its periodic time in seconds is
\begin{qoptions}
\opt{A}{$\dfrac{4\pi}{3}$}
\opt{B}{$\dfrac38\pi$}
\opt{C}{$\dfrac{8\pi}{3}$}
\opt{D}{$\dfrac73\pi$}
\end{qoptions}
\qn{18} A particle of mass $2$kg moves in simple harmonic motion and its potential energy $U$ varies with position $x$ as shown. The period of oscillation of the particle is :-
\qfig{p1_phy_s1_q18_1}
\begin{qoptions}
\opt{A}{$\dfrac{2\pi}{5}$s}
\opt{B}{$\dfrac{2\sqrt2\pi}{5}$s}
\opt{C}{$\dfrac{\sqrt2\pi}{5}$s}
\opt{D}{$\dfrac{4\pi}{5}$s}
\end{qoptions}
\qn{19} An ideal diatomic gas undergoes a process in which its internal energy $U$ and density $\rho$ vary as $U\rho=$ constant. The ratio of change in internal energy and the work done by the gas is :-
\begin{qoptions}
\opt{A}{$\dfrac32$}
\opt{B}{$\dfrac35$}
\opt{C}{$\dfrac52$}
\opt{D}{$\dfrac53$}
\end{qoptions}
\qn{20} In a three dimensional coordinate system (OXYZ), a concave mirror of radius of curvature $40$ cm is placed at $x=80$ cm. An object placed at origin is given a velocity $\vec v_0=\left(9\hat i+6\hat j+3\hat k\right)$ cm/s. Find out the magnitude of velocity of its image.
\qfig{p1_phy_s1_q20_1}
\begin{qoptions}
\opt{A}{$\sqrt{46}$ cm/s}
\opt{B}{$\sqrt6$ cm/s}
\opt{C}{$\dfrac{\sqrt{14}}{3}$ cm/s}
\opt{D}{None of these}
\end{qoptions}
\sect{SECTION-II}
\qn{1} For an ideal gas $\left(\dfrac{C_P}{C_V}=\dfrac32\right)$ undergoing a thermodynamic process work done is related to temperature change as $dW=2RdT$ for one mole. Find molar heat capacity of the gas is $xR$, then value of $3x$ is :
\qn{2} A Carnot heat engine, having efficiency $\eta=0.1$ is used as a refrigerator. If the work done on the system is $x$J, then the amount of energy absorbed from the cold reservoir is $N\,x$J, value of $N$ is
\qn{3} A uniform U-tube contains liquid mercury at same level in both limbs. Water is added in one of the limbs so that mercury level in other rises by $2$ cm, Length of water column added in the limb is `$h$' cm, then $5h$ is $[\rho_{Hg}=13.6$ g/cc, $\rho_W=1$ g/cc$]$
\qn{4} Three boats of base areas $A$, $2A$ \& $3A$, connected by a string are being pulled by an external force $F_0$. The viscosity of water is $\eta$ and depth of the water body is $H$. When the system attains a constant speed. ratio of $\dfrac{T_2}{T_1}$ is.

Given $H=5$m, $\eta=20$ poise
\qfig{p1_phy_s2_q04_1}
\qn{5} A block of mass $m$ is kept over a fixed smooth wedge. Block is attached to a sphere of same mass through fixed massless pullies $P_1$ and $P_2$. Sphere is dipped inside the water as shown. If specific gravity of material of sphere is $2$. Find the acceleration of sphere.
\qfig{p1_phy_s2_q05_1}
\qn{6} When an object is kept at a distance of $30$ cm from a concave mirror, the image is formed at a distance of $10$ cm. If the object is moved toward the mirror with a speed of $9$ m/s, the speed with which image moves is \underline{\hspace{0.8cm}} m/s.
\qn{7} A block of weight $200$ N is suspended by copper and steel wires of same cross sectional area $0.5$ cm$^2$ and, length $\dfrac{\sqrt3}{2}$m and $\dfrac12$m, respectively. Their other ends are fixed on a ceiling as shown in figure. The angles subtended by copper and steel wires with ceiling are $30^\circ$ and $60^\circ$, respectively. If elongation in copper wire is ($\Delta\ell_1$) and elongation in steel wire is ($\Delta\ell_2$), then the ratio $\dfrac{\Delta\ell_1}{\Delta\ell_2}$ is \underline{\hspace{0.8cm}}.

[Young's modulus for copper and steel are $1\times10^{11}$ N/m$^2$ and $2\times10^{11}$ N/m$^2$, respectively.]
\qfig{p1_phy_s2_q07_1}
\qn{8} A tank of cross-sectional area $A_0$ is filled with a liquid. A small orifice of area $2A$ ($A\ll\dfrac{A_0}{2}$) is present at the bottom of tank. At a moment the height of liquid in the tank is `$H$', acceleration of top layer of liquid at this moment comes out $\left(\dfrac{2A}{A_0}\right)^2k$ find value of $k$. Assume that the viscosity of the liquid is negligible. $[g=10$ m/s$^2]$
\qfig{p1_phy_s2_q08_1}
\qn{9} A tank contains three layers of immiscible liquids. The first layer is of water with refractive index $4/3$ and thickness $4$ cm. The second layer is of oil with refractive index $3/2$ and thickness $6$ cm while third layer is of liquid with refractive index $2$ and thickness $8$ cm. The shift in the position of bottom when viewed from the above as shown is \underline{\hspace{0.8cm}} cm.
\qfig{p1_phy_s2_q09_1}
\qn{10} Distance between image and object placed at origin is `$n$' cm, then $n$ is
\qfig{p1_phy_s2_q10_1}
\partsec{2}{CHEMISTRY}
\sect{SECTION-I}
\qn{1} Which of the following match is incorrect for the thermodynamically most stable allotropic form of that element?
\begin{qoptions}
\opt{A}{C$\to$Graphite}
\opt{B}{S$\to$Rhombic sulphur}
\opt{C}{P$\to$Black phosphorus}
\opt{D}{None of these}
\end{qoptions}
\qn{2} Which of the following set is not correct
\begin{qoptions}
\opt{A}{$\mathrm{SO_3}$, $\mathrm{O_3}$, $\mathrm{NH_4^+}$ all have coordinate bonds}
\opt{B}{$\mathrm{H_2O}$, $\mathrm{NO_2^-}$, $\mathrm{ClO_2^-}$ all are v shaped}
\opt{C}{$\mathrm{I_3^-}$, $\mathrm{ICl_2^-}$, $\mathrm{NO_2^+}$ all are linear species}
\opt{D}{$\mathrm{SF_4}$, $\mathrm{SiF_4}$, $\mathrm{XeF_4}$ are tetrahedral in shape}
\end{qoptions}
\qn{3} Which of the following lewis acid - lewis base interaction is NOT possible.
\begin{qoptions}
\opt{A}{$\mathrm{(CH_3)_2O+BF_3\longrightarrow(CH_3)_2O.BF_3}$}
\opt{B}{$\mathrm{BeF_2+4F^-\longrightarrow BeF_6^{4-}}$}
\opt{C}{$\mathrm{PCl_5+Cl^-\longrightarrow PCl_6^-}$}
\opt{D}{$\mathrm{SiF_4+2F^-\longrightarrow SiF_6^{2-}}$}
\end{qoptions}
\qn{4} Specie containing maximum number of identical bond angles among the following is.
\begin{qoptions}
\opt{A}{$\mathrm{CH_2F_2}$}
\opt{B}{$\mathrm{XeOF_4}$}
\opt{C}{$\mathrm{PCl_3F_2}$}
\opt{D}{$\mathrm{XeF_5^{\ominus}}$}
\end{qoptions}
\qn{5} Among the following which does not have same value of dipole moment as $\mathrm{C_6H_5Cl}$.
\begin{qoptions}
\opt{A}{\qfigi{p1_chem_s1_q05_A}}
\opt{B}{\qfigi{p1_chem_s1_q05_B}}
\opt{C}{m-$\mathrm{C_6H_4Cl_2}$}
\opt{D}{p-$\mathrm{C_6H_4Cl_2}$}
\end{qoptions}
\qn{6} Which of the following statement is \textbf{INCORRECT}?
\begin{qoptions}
\opt{A}{HF has higher boiling point than $\mathrm{NH_3}$}
\opt{B}{$\mathrm{CO_2}$, $\mathrm{N_3^-}$ and $\mathrm{NO_2^+}$ are linear}
\opt{C}{Order of $\sigma$-bond strength $2s$-$2s<2s$-$2p<2p$-$2p$}
\opt{D}{Among $\mathrm{CH_3^+}$, $\mathrm{AlCl_3}$, $\mathrm{AlF_3}$ and $\mathrm{BF_3}$, only $\mathrm{AlF_3}$ is hypovalent species.}
\end{qoptions}
\qn{7} de-Broglie wavelength associated with He-atom is directly proportional to
\begin{qoptions}
\opt{A}{Molar mass}
\opt{B}{$\dfrac{1}{\text{Molar mass}}$}
\opt{C}{$\dfrac{1}{\sqrt{\text{Temperature}}}$}
\opt{D}{Temperature}
\end{qoptions}
\qn{8} For the reaction :

$\mathrm{K_4[Fe(CN)_6]\longrightarrow Fe^{3+}+CO_2+NO_3^-}$

The equivalent mass of $\mathrm{K_4[Fe(CN)_6]}$ is
\begin{qoptions}
\opt{A}{$\dfrac53\times$ Molar mass}
\opt{B}{$\dfrac{1}{61}\times$ Molar mass}
\opt{C}{Molar mass}
\opt{D}{$11\times$ Molar mass}
\end{qoptions}
\qn{9} Consider the ionisation

$\mathrm{H^+(aq)+H_3N^+CH_2COO^-(aq)\rightleftharpoons}$

$\mathrm{H_3N^+CH_2COOH(aq)\rightleftharpoons H_2NCH_2COOH(aq)+H^+(aq)}$

Assuming that the following acid dissociation constant apply to the ionisations :

{\small$\mathrm{H_3N^+CH_2COOH(aq)\rightleftharpoons H^+(aq)+H_3N^+CH_2COO^-(aq)}$}; $K_a=4\times10^{-3}$

$\mathrm{H_3N^+CH_2COO^-(aq)\rightleftharpoons H_2NCH_2COOH(aq)+H^+(aq)}$; $K_a=1.6\times10^{-10}$

Estimate a value for the equilibrium constant for the process

$\mathrm{H_3N^+CH_2COO^-(aq)\rightleftharpoons H_2NCH_2COOH(aq)}$ is
\begin{qoptions}
\opt{A}{$4\times10^{-8}$}
\opt{B}{$0.25\times10^{7}$}
\opt{C}{$6.4\times10^{-13}$}
\opt{D}{$0.25\times10^{13}$}
\end{qoptions}
\qn{10} A mixture of $\mathrm{NH_3(g)}$ and $\mathrm{N_2H_4(g)}$ is placed in sealed container at $300$K. The total pressure is $0.5$ atm. The container is heated to $1200$ K at which $\mathrm{NH_3(g)}$ and $\mathrm{N_2H_4(g)}$ decomposes completely as follows:

$\mathrm{2NH_3(g)\longrightarrow N_2(g)+3H_2(g)}$

$\mathrm{N_2H_4(g)\longrightarrow N_2(g)+2H_2(g)}$

After decomposition, the total pressure at $1200$ K is found to be $4.5$ atm. The mole\% of $\mathrm{N_2H_4}$ in the original mixture at $300$ K is :
\begin{qoptions}
\opt{A}{$20\%$}
\opt{B}{$25\%$}
\opt{C}{$50\%$}
\opt{D}{$75\%$}
\end{qoptions}
\qn{11} A reaction at equilibrium involving $2$ moles each $\mathrm{PCl_5}$, $\mathrm{PCl_3}$ and $\mathrm{Cl_2}$ is maintained at $250^\circ$C and a total pressure of $12$ atm. The value of $K_P$ for the reaction

$\dfrac12\mathrm{PCl_3}+\dfrac12\mathrm{Cl_2}\rightleftharpoons\dfrac12\mathrm{PCl_5}$
\begin{qoptions}
\opt{A}{$2$ atm$^{-1/2}$}
\opt{B}{$4$ atm$^{-1/2}$}
\opt{C}{$0.25$ atm$^{-1/2}$}
\opt{D}{$0.5$ atm$^{-1/2}$}
\end{qoptions}
\qn{12} Which of the following graph is/are inconsistent with ideal gas behaviours?
\begin{qoptions}
\opt{A}{\qfigi{p1_chem_s1_q12_A}}
\opt{B}{\qfigi{p1_chem_s1_q12_B}}
\opt{C}{\qfigi{p1_chem_s1_q12_C}}
\opt{D}{\qfigi{p1_chem_s1_q12_D}}
\end{qoptions}
\qn{13} For the reaction : $\mathrm{2SO_2(g)+O_2(g)\rightleftharpoons2SO_3(g)}$; $\Delta H=-$ve. An increase in temperature shows.
\begin{qoptions}
\opt{A}{More dissociation of $\mathrm{SO_3}$ and a decrease in $K_c$}
\opt{B}{Less dissociation of $\mathrm{SO_3}$ and an increase in $K_c$}
\opt{C}{More dissociation of $\mathrm{SO_3}$ and an increase in $K_c$}
\opt{D}{Less dissociation of $\mathrm{SO_3}$ and a decrease in $K_c$}
\end{qoptions}
\qn{14} Identify correct acidic strength order?
\begin{qoptions}
\opt{A}{\qfigi{p1_chem_s1_q14_A}}
\opt{B}{\qfigi{p1_chem_s1_q14_2}\ \qfigi{p1_chem_s1_q14_3}\ $>$\ \qfigi{p1_chem_s1_q14_4}}
\opt{C}{Me--C$\equiv$C--H $>$ MeOH $>$ $\mathrm{NH_3}$}
\opt{D}{\qfigi{p1_chem_s1_q14_D}}
\end{qoptions}
\qn{15} Choose the order that has the following carbocations correctly arranged with respect to increasing stability?
\qfig{p1_chem_s1_q15_1}
\qfig{p1_chem_s1_q15_2}
\qfig{p1_chem_s1_q15_3}
\begin{qoptions}
\opt{A}{$1>2>3$}
\opt{B}{$1>3>2$}
\opt{C}{$2>1>3$}
\opt{D}{$3>2>1$}
\end{qoptions}
\qn{16} Identify major product of following reaction?
\qfig{p1_chem_s1_q16_1}
\begin{qoptions}
\opt{A}{\qfigi{p1_chem_s1_q16_A}}
\opt{B}{\qfigi{p1_chem_s1_q16_B}}
\opt{C}{\qfigi{p1_chem_s1_q16_C}}
\opt{D}{\qfigi{p1_chem_s1_q16_D}}
\end{qoptions}
\qn{17} Identify major product of following reaction?
\qfig{p1_chem_s1_q17_1}
\begin{qoptions}
\opt{A}{\qfigi{p1_chem_s1_q17_A}}
\opt{B}{\qfigi{p1_chem_s1_q17_B}}
\opt{C}{\qfigi{p1_chem_s1_q17_C}}
\opt{D}{$\mathrm{CH_2{=}CH{-}CH{=}CH{-}CH{=}CH_2}$}
\end{qoptions}
\qn{18} Identify reaction which is not correctly matched with its major products?
\begin{qoptions}
\opt{A}{\qfigi{p1_chem_s1_q18_1}\ \qfigi{p1_chem_s1_q18_2}}
\opt{B}{\qfigi{p1_chem_s1_q18_B}}
\opt{C}{\qfigi{p1_chem_s1_q18_C}\ $\to$\ \qfigi{p1_chem_s1_q18_5}}
\opt{D}{\qfigi{p1_chem_s1_q18_D}\ \qfigi{p1_chem_s1_q18_7}}
\end{qoptions}
\qn{19} Identify Diastereoisomeric pair of compounds?
\begin{qoptions}
\opt{A}{\qfigi{p1_chem_s1_q19_A}}
\opt{B}{\qfigi{p1_chem_s1_q19_B}}
\opt{C}{\qfigi{p1_chem_s1_q19_C}}
\opt{D}{\qfigi{p1_chem_s1_q19_D}}
\end{qoptions}
\qn{20} 
\qfig{p1_chem_s1_q20_1}
\qfig{p1_chem_s1_q20_2}
Compounds are seprated by?
\begin{qoptions}
\opt{A}{Column chromatography}
\opt{B}{Filtration}
\opt{C}{Duma method}
\opt{D}{Vaccum Distillation}
\end{qoptions}
\sect{SECTION-II}
\qn{1} How many of the following have atleast one bond angle greater than HNH in $\mathrm{NH_3}$.

$\mathrm{O(CH_3)_2}$, $\mathrm{B(OH)_3}$, $\mathrm{BF_3}$, $\mathrm{XeO_3F_2}$, $\mathrm{SiH_4}$, $\mathrm{SO_2}$
\qn{2} Identify the species having ionic bond as well as p$\pi$--d$\pi$ type of bonds

$\mathrm{CsClO_4}$, $\mathrm{Na_2SO_4}$, $\mathrm{RbNO_3}$, $\mathrm{SO_3}$, $\mathrm{KHCO_3}$, $\mathrm{CaCO_3}$, $\mathrm{AlCl_3}$
\qn{3} Let the number of LP in the molecule $=x$

Let the number of $\sigma$ bonds in the molecule $=y$

Let the number of $\pi$ bonds in the molecules $=z$

Then $=\dfrac{x\times y}{z}$ in $\mathrm{NSF_3}$ is
\qn{4} Find the value of $(x+y)$ during lewis acid base interaction of given species.

$x=$ number of species in which planarity/nonplanarity around underlined atom is changed by donation of lone pair among $\mathrm{H_2\underline{O}}$, $\mathrm{\underline{N}H_3}$, $\mathrm{\underline{P}Cl_3}$

$y=$ Number of species in which planarity/Non planarity around underlined atom is changed by acceptance of one lone pair among $\mathrm{\underline{B}F_3}$, $\mathrm{\underline{P}Cl_5}$, $\mathrm{\underline{Xe}F_4}$
\qn{5} Calculate the pressure exerted (in atm) by $1$ mole of $\mathrm{CH_4}$ gas at a temperature of $\dfrac{24}{0.0821}$ K if volume occupied by $\mathrm{CH_4}$ molecules is negligible. Given $a=2$ atm lit$^2$ mol$^{-2}$.
\qn{6} In a container a gaseous mixture containing He and $\mathrm{SO_2}$ gas present in $1:4$ mole ratio allowed through an orifice. Vapour density of initially diffused mixture is \underline{\hspace{1cm}}.
\qn{7} Pure gas `A' was taken in a closed container at an initial pressure of $5$ atm. The pressure in the container changed due to reactions

$\mathrm{A(g)}\overset{K_{p_1}}{\rightleftharpoons}\mathrm{B(g)}+2\mathrm{C(g)}$

$\mathrm{C(g)}\overset{K_{p_2}}{\rightleftharpoons}\mathrm{B(g)}+\mathrm{D(g)}$

If the total pressure due to mixture of gas A, B, C and D was found to be $10$ atm, in which pressure due to C was $3$ atm. Find $K_{P_1}$ (atm$^2$) for first equilibria.
\qn{8} Identify total number of compounds which are optically active?
\qfig{p1_chem_s2_q08_1}
\qfig{p1_chem_s2_q08_2}
\qfig{p1_chem_s2_q08_3}
\qfig{p1_chem_s2_q08_4}
\qn{9} Identify total number of carbocations which undergoes rearrangement reactions
\qfig{p1_chem_s2_q09_1}
\qfig{p1_chem_s2_q09_2}
\qfig{p1_chem_s2_q09_3}
\qfig{p1_chem_s2_q09_4}
\qfig{p1_chem_s2_q09_5}
\qfig{p1_chem_s2_q09_6}
\qn{10} Total number of stereo isomers for following structure?
\qfig{p1_chem_s2_q10_1}
\partsec{3}{MATHEMATICS}
\sect{SECTION-I}
\qn{1} If $A=\{1,3,5,7,11,13,15,17,19,21\}$

$B=\{2,4,6,\ldots\ldots,22\}$

$U$ is universal set, then $A'\cup((A\cup B)\cap B')$ is -
\begin{qoptions}
\opt{A}{$U$}
\opt{B}{$A'$}
\opt{C}{$A\cap B$}
\opt{D}{$A$}
\end{qoptions}
\qn{2} The value of $\displaystyle\sum_{r=0}^{n}(a+rd)\,{}^nC_r$ (where $a$ and $d$ are constants) is equal to -
\begin{qoptions}
\opt{A}{$a.2^n$}
\opt{B}{$(a+(2n-1)d)2^{n-1}$}
\opt{C}{$(2a+nd)2^{n-1}$}
\opt{D}{$(a+nd)2^n$}
\end{qoptions}
\qn{3} If $f(x)=-2x^2+2x+20$, then the number of integer(s) in the range of $f(x)$ for $x\in[-1,1]$ is (are)
\begin{qoptions}
\opt{A}{$1$}
\opt{B}{$3$}
\opt{C}{$4$}
\opt{D}{$5$}
\end{qoptions}
\qn{4} A line inclined at an angle of $135^\circ$ with positive direction of $x$-axis passes through the point $A(3,4)$. There are two points $P$ and $Q$ on the line such that $AP=AQ=2\sqrt2$. The co-ordinates of $P$ and $Q$ are respectively -
\begin{qoptions}
\opt{A}{$(1,6)$ and $(5,-2)$}
\opt{B}{$(-1,-6)$ and $(5,2)$}
\opt{C}{$(1,6)$ and $(5,2)$}
\opt{D}{$(1,-6)$ and $(-5,2)$}
\end{qoptions}
\qn{5} The number of ways in which three distinct numbers in A.P. can be selected from $1,2,3,\ldots\ldots,24$ is -
\begin{qoptions}
\opt{A}{$132$}
\opt{B}{$572$}
\opt{C}{$264$}
\opt{D}{$150$}
\end{qoptions}
\qn{6} The least positive integral value of `$K$' for which $(x+k-1)(x-2k-7)<0\ \forall x\in[-1,2]$ is
\begin{qoptions}
\opt{A}{$4$}
\opt{B}{$3$}
\opt{C}{$2$}
\opt{D}{$1$}
\end{qoptions}
\qn{7} The solution set of equation $\log_{\cos x}\left(\dfrac{2\sin x+3}{4}\right)\geqslant2$, where $x\in[0,\pi]$, is
\begin{qoptions}
\opt{A}{$\left[0,\dfrac\pi5\right]$}
\opt{B}{$\left(0,\dfrac{\pi}{10}\right]$}
\opt{C}{$\left[0,\dfrac{\pi}{10}\right]$}
\opt{D}{$\phi$ (null set)}
\end{qoptions}
\qn{8} The number of ways in which the letters of the word `LONDON' can be rearranged if the two `O' s are together but the two `N's are seperated is
\begin{qoptions}
\opt{A}{$72$}
\opt{B}{$71$}
\opt{C}{$36$}
\opt{D}{$35$}
\end{qoptions}
\qn{9} If $(201)!$ is divided by $24^k$, then the largest value of $k$ is
\begin{qoptions}
\opt{A}{$98$}
\opt{B}{$65$}
\opt{C}{$49$}
\opt{D}{$66$}
\end{qoptions}
\qn{10} If $ax+by=10$ is the chord of minimum length of the circle $(x-10)^2+(y-20)^2=729$ and the chord passes through $(5,15)$, then the value of $(4a+2b)$ is
\begin{qoptions}
\opt{A}{$2$}
\opt{B}{$3$}
\opt{C}{$4$}
\opt{D}{$5$}
\end{qoptions}
\qn{11} If $y=2x+k$ is diameter to the circle $2(x^2+y^2)+3x+4y-1=0$, then $k$ is
\begin{qoptions}
\opt{A}{$0$}
\opt{B}{$\dfrac12$}
\opt{C}{$1$}
\opt{D}{$2$}
\end{qoptions}
\qn{12} Consider the circle $S:x^2+y^2=1$ and the line $L:y-m(x+2)=0$. If $L$ intersects $S$ at $P$ and $Q$, then locus of middle point of $PQ$ is
\begin{qoptions}
\opt{A}{$(x+1)^2+y^2=1$}
\opt{B}{$x^2+(y-1)^2=1$}
\opt{C}{$(x-1)^2+y^2=1$}
\opt{D}{$x^2+(y+1)^2=1$}
\end{qoptions}
\qn{13} The number of ways of selection of $4$ letters of the word `PERFECT ' is
\begin{qoptions}
\opt{A}{$15$}
\opt{B}{$21$}
\opt{C}{$25$}
\opt{D}{$31$}
\end{qoptions}
\qn{14} The number of points $(x,y)$ having integral coordinates satisfying $x^2+y^2<25$ is -
\begin{qoptions}
\opt{A}{$69$}
\opt{B}{$80$}
\opt{C}{$81$}
\opt{D}{$77$}
\end{qoptions}
\qn{15} In triangle ABC, if $A-B=90^\circ$ and $R=4r$, where $R$ and $r$ have their usual meanings, then $\cos C$ is equal to
\begin{qoptions}
\opt{A}{$\dfrac23$}
\opt{B}{$\dfrac78$}
\opt{C}{$\dfrac34$}
\opt{D}{$\dfrac{\sqrt3}{2}$}
\end{qoptions}
\qn{16} If the line $y=x$ and $x+y=0$ are tangents to a variable circle, then locus of the centre of the variable circle, is
\begin{qoptions}
\opt{A}{$xy=0$}
\opt{B}{$x+2y=0$}
\opt{C}{$2x+3y=5$}
\opt{D}{$2x-3y=5$}
\end{qoptions}
\qn{17} The sum of all $4$ digit numbers that can be formed by using the digits $2,4,6,8$. (repetition of digits not being allowed) is -
\begin{qoptions}
\opt{A}{$133320$}
\opt{B}{$123330$}
\opt{C}{$113230$}
\opt{D}{$323430$}
\end{qoptions}
\qn{18} A line drawn from the origin touches the curve $y^2-x^2=1-x$ in the first quadrant. If the line is inclined at an acute angle $\theta$ with the positive $x$ - axis , then $\dfrac{\tan\theta}{\sqrt3}$ is equal to
\begin{qoptions}
\opt{A}{$\dfrac{\sqrt3}{2}$}
\opt{B}{$\dfrac12$}
\opt{C}{$\dfrac{2}{\sqrt3}$}
\opt{D}{$\dfrac13$}
\end{qoptions}
\qn{19} If $\log_{27}x+\log_9y=\dfrac72$, $\log_{27}y+\log_9x=\dfrac23$, then the value of $xy$ is
\begin{qoptions}
\opt{A}{$27$}
\opt{B}{$81$}
\opt{C}{$243$}
\opt{D}{$216$}
\end{qoptions}
\qn{20} Consider the number $N=2^4\,3^4\,5^2\,7^2\,11^2$, then the sum of all the odd divisors which are divisible by `5' but not by `7' is $11^2(133\lambda)$ then $\lambda$ is equal to
\begin{qoptions}
\opt{A}{$10$}
\opt{B}{$30$}
\opt{C}{$50$}
\opt{D}{$330$}
\end{qoptions}
\sect{SECTION-II}
\qn{1} If larger root of $(2015x)^2-(2014)(2016)x-1=0$ is $\alpha$ and smaller root of $x^2+2014x-2015=0$ is $\beta$, then value of $\sqrt{\dfrac{\alpha-\beta}{224}}$ is
\qn{2} If the number of integers lying in the domain of function $f(x)=\sqrt{x^2-5x+6}+\sqrt{2x+8-x^2}$ is $m$ and the number of non-negative integers lying in the range of $g(x)=\log_e\left(\sqrt{4-x^2}\right)$ is $n$, then $(m+n)$ is equal to
\qn{3} If $4^x-2^{x+2}+5+\big||b-1|-3\big|=|\sin y|$, $x,y,b\in R$ then sum of the possible values of $b$, is
\qn{4} If range of $f(x)=\dfrac{3\sin^2x+3\sin x+4}{\sin^2x+\sin x+1}$ is $[p,q]$, then value of $6p-3q$ equals
\qn{5} The ratio of the length of the sides BC and AC of $\triangle$ABC to the radius of circumscribed circle are equal to $2$ and $\dfrac32$ respectively. If the ratio of the length of the bisectors of the interior angles B and C is $\dfrac{\alpha(\sqrt\alpha-1)}{\beta\sqrt\gamma}$ (where $\alpha,\beta,\gamma\in N$), then $\alpha$ is
\qn{6} If $P_n$ denotes the product of all the coefficient of $(1+x)^n$ and $9!\,P_{n+1}=10^9\,P_n$, then $n$ is equal to
\qn{7} The number of integral values of $\lambda$ for which the variable line $3x+4y-\lambda=0$ lies between the circles $x^2+y^2-2x-2y+1=0$ and $x^2+y^2-18x-2y+78=0$ without having any common points with any circle is equal to $2k$, then value of $k$ is
\qn{8} If $1.({}^{20}C_0)^2+2.({}^{20}C_2)^2+3.({}^{20}C_4)^2+\ldots.+11.({}^{20}C_{20})^2$ is equal to $3\left[{}^{8a}C_{4a}+{}^{4a}C_{2a}\right]$, then $a$ is equal to
\qn{9} If points A and B trisect intercept of line $x+y=8$ between coordinate axes and O is origin. then value of the sum of slopes of sides of $\triangle$OAB is $k$ then greatest integer less than or equal to $k$ is
\qn{10} Let $x$ is the number of $6$ digit even numbers such that sum of their digits is $50$, then $\dfrac{x}{10}$ is equal to
\raggedcolumns
\papersec{2}
\partsec{1}{PHYSICS}
\sect{SECTION-I(i)}
\qn{1} On a smooth horizontal surface there is a disc of mass $2m$ and radius $R$ is placed. An insect of mass $m$ sits gently at a distance $\dfrac R2$ from center of disc and crawls uniformly with speed `$v_0$' with respect to centre of disc along circular path with respect to centre of disc.
\begin{qoptions}
\opt{A}{With respect to a ground observer insect moves in circular path of radius $\dfrac R2$.}
\opt{B}{With respect to a ground observer centre of disc moves on a circular path of radius $\dfrac R6$.}
\opt{C}{Net external force on this system is zero.}
\opt{D}{Insect applies a force $\dfrac{4}{27}mv_0^2R$ on disc tangentially to the surface of disc.}
\end{qoptions}
\qn{2} In a water pistol, a piston drives water through a large tube of area $A_1$ into a smaller tube of area $A_2$ as shown in figure. The radius of the large tube is $1.00$ cm and that of the small tube is $1.00$ mm. The smaller tube is $3.00$ cm above the larger tube. If the pistol is fired horizontally at a height of $1.50$ m, assume atmospheric pressure is $1.00$ atm. If the desired range of the stream is $8.00$ m,
\qfig{p2_phy_s1_q02_1}
\begin{qoptions}
\opt{A}{The speed $v_1$ with which the plunger moves is $0.146$ m/s.}
\opt{B}{The pressure inside the larger tube is $2\times10^5$ Pa}
\opt{C}{The pressure inside the larger tube is $3\times10^5$ Pa.}
\opt{D}{The force exerted on the trigger to achieve the desired range is $66$ N.}
\end{qoptions}
\qn{3} An ideal monoatomic gas is taken through a thermodynamics process $P=aV+b$ when $P$ and $V$ and pressure and volume of gas and $a$ and $b$ are positive constants.
\begin{qoptions}
\opt{A}{If $b=0$, molar heat capacity of gas is constant.}
\opt{B}{If $a\neq0$, $b\neq0$, molar heat capacity of gas is lower in higher volume}
\opt{C}{If $a=0$, increase in volume per unit rise in temperature is constant.}
\opt{D}{Slope of temperature versus volume graph is constant.}
\end{qoptions}
\qn{4} Figure is an overhead view of a mirror maze based on floor sections that are equilateral triangles. Every wall within the maze is mirrored. If you stand at entrance $x$, which of the maze monsters $a$, $b$ and $c$ hiding in the maze can you see along the virtual hall ways extending from entrance $x$ ;
\qfig{p2_phy_s1_q04_1}
\begin{qoptions}
\opt{A}{a}
\opt{B}{b}
\opt{C}{c}
\opt{D}{None}
\end{qoptions}
\qn{5} Figure shows three leg table, with top as a uniform equilateral triangular board of mass $10$ kg. Two bodies of masses $2$ kg and $3$ kg are placed at corners A and C. The table is placed on a horizontal surface.
\qfig{p2_phy_s1_q05_1}
\begin{qoptions}
\opt{A}{Normal force on leg at A $=\dfrac{160}{3}$N}
\opt{B}{Normal force on leg at B $=\dfrac{100}{3}$N}
\opt{C}{Normal force on leg at C $=\dfrac{190}{3}$N}
\opt{D}{If both the blocks are removed, normal force on leg at B will remain same.}
\end{qoptions}
\qn{6} A fish watcher at point $P$ watches a fish through a glass wall of a fish tank. The watcher is level with the fish ; the index of refraction of the glass is $8/5$, and that of the water is $4/3$. The distances are $d_1=8.0$ cm, $d_2=3.0$ cm, and $d_3=6.8$ cm.
\qfig{p2_phy_s1_q06_1}
\begin{qoptions}
\opt{A}{The apparent separation between watcher and fish as observed by the watcher is $20$ cm.}
\opt{B}{The apparent separation between watcher and fish as observed by the watcher is $13.2$ cm.}
\opt{C}{The apparent separation between watcher and fish as observed by the fish is $15$ cm.}
\opt{D}{The apparent separation between watcher and fish as observed by the fish is $7$ cm.}
\end{qoptions}
\qn{7} $20$ gm ice is kept in a calorimeter at $0^\circ$C. Now $40$ gm water at $30^\circ$C is poured gently into it. There is no leakage of heat. In equilibrium.
\begin{qoptions}
\opt{A}{Final temperature will be $0^\circ$C only if heat capacity of calorimeter is neglected.}
\opt{B}{Final temperature will be $0^\circ$C irrespective of value of heat capacity of calorimeter.}
\opt{C}{$75\%$ of ice will melt.}
\opt{D}{Remaining ice will be more than $10$ gm.}
\end{qoptions}
\qn{8} Imagine a point source $s$ which is a black body at a constant temperature $T$ at the centre of a thick spherical shell of inner radius $a$ and external radius $b$. The coefficient of thermal conductivity of the material of the shell is $K$. The inner surface of the shell is coated with lamp black. $A$ is the surface area of sources.
\qfig{p2_phy_s1_q08_1}
\begin{qoptions}
\opt{A}{The amount of thermal radiation received by the inner surface from $s$ is ($\sigma\to$ Stefan's constant) $A\sigma T^4$.}
\opt{B}{The temperature difference across the inner and outer surface is $\dfrac{A\sigma T^4}{4\pi K}\left[\dfrac1a-\dfrac1b\right]$.}
\opt{C}{If the external surface of the shell is in vacuum at $0$ K (say) the temperature of inner surface, is $\dfrac{A\sigma T^4}{4\pi K}\left[\dfrac1a-\dfrac1b\right]$.}
\opt{D}{None of these}
\end{qoptions}
\sect{SECTION-I(ii)}
\qn{9} Figure shows a rider performing a stunt with rear wheel in pure rolling.

$M=$ total mass of system $=200$ kg,

$m=$ mass of rear wheel (ring) $=10$ kg,

$R=$ radius of wheel $=50$ cm,

$h=$ height of centre of mass of system,

$b=$ distance of vertical line through centre of mass from contact point of rear wheel.
\qfig{p2_phy_s2_q09_1}
\begin{center}
{\footnotesize\setlength{\tabcolsep}{2pt}\begin{tabular}{|c|p{4.2cm}|c|p{1.4cm}|}\hline
 & \multicolumn{1}{c|}{List-I} & & List-II\\\hline
(I) & Frictional force in KN & (P) & $1.2$\\\hline
(II) & Torque provided by engine in KN-m & (Q) & $0.630$\\\hline
(III) & Value of $h$ in meter & (R) & $50$\\\hline
(IV) & Value of $b$ in centimeter & (S) & $1$\\\hline
 & & (T) & $60$\\\hline
 & & (U) & $0.525$\\\hline
\end{tabular}}
\end{center}
It the motorcycle in the given situation accelerates with $5$ m/s$^2$, the correct matching is as :-
\begin{qoptions}
\opt{A}{I $\to$ S;II $\to$ U;III $\to$ S;IV $\to$ R}
\opt{B}{I $\to$ S;II $\to$ U;III $\to$ R;IV $\to$ S}
\opt{C}{I $\to$ S;II $\to$ Q;III $\to$ S;IV $\to$ S}
\opt{D}{I $\to$ P;II $\to$ Q;III $\to$ R;IV $\to$ S}
\end{qoptions}
\qn{10} Figure shows a rider performing a stunt with rear wheel in pure rolling.

$M=$ total man of system $=200$ kg,

$m=$ mass of rear wheel (ring) $=10$ kg,

$R=$ radius of wheel $=50$ cm,

$h=$ height of centre of mass of system,

$b=$ distance of vertical line through centre of mass from contact point of rear wheel.
\qfig{p2_phy_s2_q10_1}
\begin{center}
{\footnotesize\setlength{\tabcolsep}{2pt}\begin{tabular}{|c|p{4.2cm}|c|p{1.4cm}|}\hline
 & \multicolumn{1}{c|}{List-I} & & List-II\\\hline
(I) & Frictional force in KN & (P) & $1.2$\\\hline
(II) & Torque provided by engine in KN-m & (Q) & $0.630$\\\hline
(III) & Value of $h$ in meter & (R) & $50$\\\hline
(IV) & Value of $b$ in centimeter & (S) & $1$\\\hline
 & & (T) & $60$\\\hline
 & & (U) & $0.525$\\\hline
\end{tabular}}
\end{center}
It the motorcycle in the given situation accelerates with $6$ m/s$^2$, the correct matching is as :-
\begin{qoptions}
\opt{A}{I $\to$ P;II $\to$ Q;III $\to$ S;IV $\to$ T}
\opt{B}{I $\to$ P;II $\to$ Q;III $\to$ T;IV $\to$ S}
\opt{C}{I $\to$ P;II $\to$ S;III $\to$ R;IV $\to$ Q}
\opt{D}{I $\to$ R;II $\to$ S;III $\to$ T;IV $\to$ U}
\end{qoptions}
\qn{11} A mass of $1$kg is supported by a cord $5$m long attached to a massless ring free to move on a horizontal fixed rod, as shown in the figure. Coefficient of friction between ring and the rod is $0.75$. A second cord is attached to the weight and passes through a pulley at $\dfrac{25}{3}$m from ring. Block A of mass $m$ is attached to the other end of the cord. Just after released.
\qfig{p2_phy_s2_q11_1}
\begin{center}
{\footnotesize\setlength{\tabcolsep}{2pt}\begin{tabular}{|c|p{4.2cm}|c|p{1.4cm}|}\hline
 & \multicolumn{1}{c|}{List-I} & & List-II\\\hline
(I) & Tension in the string connected to block A in `N'. & (P) & $6$\\\hline
(II) & Tension in the $5$m cord in `N' & (Q) & $8$\\\hline
(III) & The value of $\theta$ is $\dfrac\pi k$ radian, possible value of $k$. & (R) & $2$\\\hline
(IV) & Acceleration of block A in m/s$^2$ & (S) & $0$\\\hline
 & & (T) & $4$\\\hline
 & & (U) & $12$\\\hline
\end{tabular}}
\end{center}
If the ring is just about to slip
\begin{qoptions}
\opt{A}{I $\to$ P;II $\to$ Q;III $\to$ R;IV $\to$ S}
\opt{B}{I $\to$ Q;II $\to$ P;III $\to$ T;IV $\to$ S}
\opt{C}{I $\to$ T;II $\to$ Q;III $\to$ R;IV $\to$ S}
\opt{D}{I $\to$ U;II $\to$ Q;III $\to$ R;IV $\to$ S}
\end{qoptions}
\qn{12} A mass of $1$kg is supported by a cord $5$m long attached to a massless ring free to move on a horizontal fixed rod, as shown in the figure. Coefficient of friction between ring and the rod is $0.75$. A second cord is attached to the weight and passes through a pulley at $\dfrac{25}{3}$m from ring. Block A of mass $m$ is attached to the other end of this cord. Just after released.
\qfig{p2_phy_s2_q12_1}
\begin{center}
{\footnotesize\setlength{\tabcolsep}{2pt}\begin{tabular}{|c|p{4.2cm}|c|p{1.4cm}|}\hline
 & \multicolumn{1}{c|}{List-I} & & List-II\\\hline
(I) & Tension in the string connected to block A in `N'. & (P) & $6$\\\hline
(II) & Tension in the $5$m cord in `N' & (Q) & $8$\\\hline
(III) & The value of $\theta$ is $\dfrac\pi k$ radian, possible value of $k$. & (R) & $2$\\\hline
(IV) & Acceleration of block A in m/s$^2$ & (S) & $0$\\\hline
 & & (T) & $4$\\\hline
 & & (U) & $12$\\\hline
\end{tabular}}
\end{center}
The value of $m$ was initially chosen such that the ring just began to slide. Now the ring is fixed to the rod and mass of block A is increased to $5m$, holding the $1$kg mass. Just after the system is released, then
\begin{qoptions}
\opt{A}{I $\to$ P;II $\to$ Q;III $\to$ R;IV $\to$ T}
\opt{B}{I $\to$ Q;II $\to$ P;III $\to$ T;IV $\to$ U}
\opt{C}{I $\to$ T;II $\to$ U;III $\to$ R;IV $\to$ Q}
\opt{D}{I $\to$ U;II $\to$ Q;III $\to$ R;IV $\to$ P}
\end{qoptions}
\sect{SECTION-II}
\qn{1} A small spherical ball of density equal to density of liquid is released from rest in a vessel filled completely with a liquid and accelerating with acceleration $5\hat i+5\hat j$ m/s$^2$. Find initial acceleration of the ball just after the release with respect to vessel.
\qfig{p2_phy_s3_q01_1}
\qn{2} One knee smooth U-shaped tube with a circular inner section area $S$ vertically, and the other is inclined to the horizon at an angle $\alpha$. A liquid of density $\rho$ is poured into the tube and weight $M$ so that its level in the inclined knee is higher than in the vertical one, which is closed by a light piston connected to a vertical spring by stiffness $k$ (see picture). Find the period of small oscillations of this system. The free fall acceleration is $g$. (Take $M=10$ kg, $k=100$ N/m, $\rho_2=10^3$ kg/m, $g=10$ m/s$^2$, $s=10^{-2}$m$^2$, $\alpha=30^\circ$.)
\qfig{p2_phy_s3_q02_1}
\qn{3} There are two small objects on an $L=3$m long wooden board, supported at its middle by a wedge. The masses of the objects are $m_1=0.2$ kg and $m_2=0.3$ kg. Between the objects there is a compressed small spring of negligible mass. The spring is compressed to a length of $\ell=0.3$ m, and tied with a thread. The wedge is exactly below the center of mass of the system. The thread is burnt, and the spring launches the object with mass $m_1$ at a speed of $v_1=2$ m/s. The time (in sec) after which the board will tilt
\qfig{p2_phy_s3_q03_1}
\qn{4} Find the efficiency (in percentage) of the cycle shown in the diagram $p$, $V$ on figure. The working fluid is one mole of an ideal monatomic gas.
\qfig{p2_phy_s3_q04_1}
\qn{5} A fluid system, contained in a piston and cylinder machine, passes through a complete cycle of four processes. The sum of all heat transferred during a cycle is $-340$ kJ. The system completes $200$ cycles per min. The following table shows the heat, work and internal energy for each process. The rate of work output (in kW) is
\begin{center}
{\footnotesize\setlength{\tabcolsep}{2pt}\begin{tabular}{|l|l|l|l|}\hline
Process & $Q$ (kJ/min) & $W$ (kJ/min) & $\Delta E$ (kJ/min)\\\hline
1-2 & $0$ & $4340$ & $-$\\\hline
2-3 & $42000$ & $0$ & $-$\\\hline
3-4 & $-4200$ & $-$ & $-73200$\\\hline
4-1 & $-$ & $-$ & $-$\\\hline
\end{tabular}}
\end{center}
\qn{6} Determine the entropy change of $4$ kg of a perfect gas whose temperature varies from $127^\circ$C to $227^\circ$C during a constant volume process. The specific heat varies linearly with absolute temperature and is represented by the relation : $c_v=(0.48+0.0096\,T)$ kJ/kg K. [Take $\ell$n $(1.25)=0.22$]
\partsec{2}{CHEMISTRY}
\sect{SECTION-I(i)}
\qn{1} In a closed rigid $2$ L container the following reaction is occuring :

$\mathrm{PCl_5(g)\rightleftharpoons PCl_3(g)+Cl_2(g)}$

If initially taken $2$ moles of $\mathrm{PCl_5}$ dissociates to $60\%$ to establish equilibrium then which of the following is/are correct statement(s) :
\begin{qoptions}
\opt{A}{Mass inside the container remains same.}
\opt{B}{Moles inside the container increases.}
\opt{C}{Molar mass of equilibrium mixture is lesser than molar mass of $\mathrm{PCl_5}$.}
\opt{D}{$K_C<1$}
\end{qoptions}
\qn{2} One of the major use of $\mathrm{H_2O_2}$ is for artificial respiration. It is stored in cylinder of $3$ litre capacity completely filled with $\mathrm{H_2O_2}$ solution of `$45.4$V'. In a hospital of kota, patient under artificial respiration takes $200$ ml $\mathrm{O_2}$ per min at $1$ atm at $273$ K for $2.8$ hours from cylinder.

[Assume volume of solution and rate of decomposition remain constant]

Then which of the following statement(s) is/are correct.
\begin{qoptions}
\opt{A}{Initial moles of $\mathrm{H_2O_2}$ in the cylinder are $12$.}
\opt{B}{Volume of oxygen used for respiration is $33.6$ L.}
\opt{C}{Volume strength of $\mathrm{H_2O_2}$ left in solution is $11.35$ V.}
\opt{D}{Volume strength of $\mathrm{H_2O_2}$ left in solution is $34.05$ V.}
\end{qoptions}
\qn{3} Identify the structure of transdecaline
\begin{qoptions}
\opt{A}{\qfigi{p2_chem_s1_q03_A}}
\opt{B}{\qfigi{p2_chem_s1_q03_B}}
\opt{C}{\qfigi{p2_chem_s1_q03_C}}
\opt{D}{\qfigi{p2_chem_s1_q03_D}}
\end{qoptions}
\qn{4} Identify the structure of identical compounds
\begin{qoptions}
\opt{A}{\qfigi{p2_chem_s1_q04_A}}
\opt{B}{\qfigi{p2_chem_s1_q04_B}}
\opt{C}{\qfigi{p2_chem_s1_q04_C}}
\opt{D}{\qfigi{p2_chem_s1_q04_D}}
\end{qoptions}
\qn{5} Identify the geometrical isomers
\begin{qoptions}
\opt{A}{\qfigi{p2_chem_s1_q05_A}}
\opt{B}{\qfigi{p2_chem_s1_q05_B}}
\opt{C}{\qfigi{p2_chem_s1_q05_C}}
\opt{D}{\qfigi{p2_chem_s1_q05_D}}
\end{qoptions}
\qn{6} Identify correct statement?
\begin{qoptions}
\opt{A}{$\mathrm{F^{\ominus}<Cl^{\ominus}<Br^{\ominus}<I^{\ominus}}$ Nucleophilicity order in $\mathrm{H_2O}$}
\opt{B}{\qfigi{p2_chem_s1_q06_B}}
\opt{C}{\qfigi{p2_chem_s1_q06_C}}
\opt{D}{\qfigi{p2_chem_s1_q06_D}}
\end{qoptions}
\qn{7} Select the correct option(s).
\begin{qoptions}
\opt{A}{BP of $\mathrm{NH_3}$ is greater than $\mathrm{SbH_3}$ due to strong H-bonding}
\opt{B}{If osmotic pressure of $\mathrm{CH_3COOH}$ solution in benzene is directly proportional to effective number of molecular units in $\mathrm{CH_3COOH}$ molecules, then the pressure is found twice than expected.}
\opt{C}{If a drop of ethanol, ethylene glycol and glycerol is poured down a wooden cardboard through their droppers, the drop of ethanol would reach the bottom first}
\opt{D}{The intermolecular spaces between water molecules decreases from below $0^\circ$C to $4^\circ$C and then increases on further increasing the temperature}
\end{qoptions}
\qn{8} Select the incorrect option(s).
\begin{qoptions}
\opt{A}{Number of P--OH and P--O--P bonds are same in $\mathrm{H_6P_4O_{13}}$ and $\mathrm{H_4P_4O_{12}}$ but number of P $=$ O bonds are different}
\opt{B}{Thiosulphuric acid is a member of thionic acid series}
\opt{C}{Difference in acidic strength of nitric acid and peroxynitric acid can be judged by observing stability through resonance in their conjugate bases}
\opt{D}{Sodium hydrogen pyrophosphite will release two ions upon dissolution in water}
\end{qoptions}
\sect{SECTION-I(ii)}
\qn{9} 
\begin{center}
{\footnotesize\setlength{\tabcolsep}{2pt}\begin{tabular}{|c|p{3.4cm}|c|p{3.4cm}|}\hline
 & \centering List-I\\(condition of vander walls gas) & & \centering List-II\\(compressibility factor)\tabularnewline\hline
(I) & At low to moderate pressure & (P) & $Z=1+\dfrac{Pb}{RT}$\\\hline
(II) & At critical condition of P, V, T & (Q) & $Z=1-\dfrac{a}{V_mRT}$\\\hline
(III) & For $\mathrm{H_2}$ \& helium gas at room temperature & (R) & $Z=3/8$\\\hline
(IV) & At very low pressure \& high temperature & (S) & $Z=1$\\\hline
 & & (T) & Positive derivation from ideality\\\hline
 & & (U) & Negative derivation from ideality\\\hline
\end{tabular}}
\end{center}
Select the CORRECT option -
\begin{qoptions}
\opt{A}{III $\to$ Q}
\opt{B}{II $\to$ S}
\opt{C}{IV $\to$ P}
\opt{D}{III $\to$ P}
\end{qoptions}
\qn{10} 
\begin{center}
{\footnotesize\setlength{\tabcolsep}{2pt}\begin{tabular}{|c|p{3.4cm}|c|p{3.4cm}|}\hline
 & \centering List-I\\(condition of vander walls gas) & & \centering List-II\\(compressibility factor)\tabularnewline\hline
(I) & At low to moderate pressure & (P) & $Z=1+\dfrac{Pb}{RT}$\\\hline
(II) & At critical condition of P, V, T & (Q) & $Z=1-\dfrac{a}{V_mRT}$\\\hline
(III) & For $\mathrm{H_2}$ \& helium gas at room temperature & (R) & $Z=3/8$\\\hline
(IV) & At very low pressure \& high temperature & (S) & $Z=1$\\\hline
 & & (T) & Positive derivation from ideality\\\hline
 & & (U) & Negative derivation from ideality\\\hline
\end{tabular}}
\end{center}
Select the CORRECT option -
\begin{qoptions}
\opt{A}{I $\to$ P}
\opt{B}{II $\to$ R}
\opt{C}{IV $\to$ T}
\opt{D}{IV $\to$ U}
\end{qoptions}
\qn{11} 
\begin{center}
{\footnotesize\setlength{\tabcolsep}{2pt}\begin{tabular}{|c|p{1.6cm}|c|p{4.2cm}|}\hline
 & \centering List-I\\(Molecules) & & \centering List-II\\(Various facts)\tabularnewline\hline
(I) & Neo--$\mathrm{C_5H_{12}}$ & (P) & Maximum number of same atoms in one plane $=3$\\\hline
(II) & $\mathrm{C_2H_2Cl_2}$ & (Q) & n-isomer has more surface area hence more BP then iso -- and neo --\\\hline
(III) & $\mathrm{CH_3Cl}$ & (R) & Cis-form has more dipole moment than trans-form\\\hline
(IV) & $\mathrm{CH_4}$ & (S) & has more boiling point than tetrafluoromethane due to more repulsive forces in latter\\\hline
 & & (T) & has more dipole moment than dichloromethane and trichloromethane\\\hline
 & & (U) & have regular geometry around carbon atom(s)\\\hline
\end{tabular}}
\end{center}
Select the incorrect match.
\begin{qoptions}
\opt{A}{I $\to$ P}
\opt{B}{II $\to$ R}
\opt{C}{III $\to$ T}
\opt{D}{IV $\to$ S}
\end{qoptions}
\qn{12} 
\begin{center}
{\footnotesize\setlength{\tabcolsep}{2pt}\begin{tabular}{|c|p{1.6cm}|c|p{4.2cm}|}\hline
 & \centering List-I\\(Molecules) & & \centering List-II\\(Various facts)\tabularnewline\hline
(I) & Neo--$\mathrm{C_5H_{12}}$ & (P) & Maximum number of same atoms in one plane $=3$\\\hline
(II) & $\mathrm{C_2H_2Cl_2}$ & (Q) & n-isomer has more surface area hence more BP then iso -- and neo --\\\hline
(III) & $\mathrm{CH_3Cl}$ & (R) & Cis-form has more dipole moment than trans-form\\\hline
(IV) & $\mathrm{CH_4}$ & (S) & has more boiling point than tetrafluoromethane due to more repulsive forces in latter\\\hline
 & & (T) & has more dipole moment than dichloromethane and trichloromethane\\\hline
 & & (U) & have regular geometry around carbon atom(s)\\\hline
\end{tabular}}
\end{center}
Select the correct match.
\begin{qoptions}
\opt{A}{I $\to$ U}
\opt{B}{II $\to$ P}
\opt{C}{III $\to$ P}
\opt{D}{IV $\to$ T}
\end{qoptions}
\sect{SECTION-II}
\qn{1} $10$ g impure mixture of $\mathrm{KIO_3}$ was titrated with excess of KI in acidic medium. Iodine so liberated required $100$ ml of $0.15$ N solution of sodium thiosulphate for titration. The mass-percentage of $\mathrm{KIO_3}$ in the mixture is \underline{\hspace{1cm}}.

(Atomic mass : K $=39$, I $=127$, O $=16$)
\qn{2} The wave motion of electron in a Bohr's orbit of hydrogen is as shown in diagram. The potential energy of electron in the given orbit of hydrogen atoms $\dfrac{(-13.6)\times2}{x}$. Then value of $x$ is :
\qfig{p2_chem_s3_q02_1}
\qn{3} In carius method of estimation of halogen $0.15$ g of an organic compound gave $0.12$ g of AgBr. Find out the percentage of bromine in the compound?
\qn{4} On complete combustion, $0.246$ g of an organic compound gave $0.198$ g of carbon dioxide and $0.1014$g of water. Determine the percentage composition of carbon in the compound
\qn{5} Consider the following molecules

$\mathrm{CO_2}$, $\mathrm{BF_3}$, $\mathrm{H_2O}$, $\mathrm{NH_3}$, $\mathrm{CH_4}$, $\mathrm{PCl_5}$, $\mathrm{SF_6}$, $\mathrm{IF_7}$, $\mathrm{PCl_3F_2}$, $\mathrm{PCl_2F_3}$, $\mathrm{XeF_5^{\ominus}}$, $\mathrm{ClF_3}$, $\mathrm{SF_4}$

$\dfrac{\text{number of molecules having all equal bond length}}{\text{number of polar molecules showing H - bonding}}=x$

Find $x$?
\qn{6} Regarding lewis dot structure of sulphate ion, consider the following statements

(a) Resonance is possible with bond order $=1.5$

(b) Each atom contains formal charge

(c) The expanded structure and lewis structure have equal number of S--O linkage

(d) Magnitude of formal charge on sulphur is equal to $\pi$ bonds in most stable expanded structure

(e) Ratio of total number of lone pairs to bond pairs is $13:4$

(f) Formal charge of sulphur here is different from formal charge of sulphur in lewis dot structure of $\mathrm{SO_3}$ and $\mathrm{SO_2Cl_2}$

$\dfrac{\text{number of correct statements}}{\text{number of wrong statements}}=x$

Find $x$
\partsec{3}{MATHEMATICS}
\sect{SECTION-I(i)}
\qn{1} Consider the curve

$S=(\sin\theta)x^2+2(\cos\theta)xy+(\sin\theta)y^2+2(\cos\theta)x+2(\cos\theta)y+\sin\theta=0$

Now, choose the correct statement(s)
\begin{qoptions}
\opt{A}{If $\theta=\dfrac{3\pi}{4}$, then $S=0$ represents a pair of straight lines.}
\opt{B}{If $S=0$ represents a pair of real straight lines, then $\theta=n\pi+\dfrac\pi4$, $n\in Z$}
\opt{C}{If $S=0$ represents a pair of real straight lines, then they are parallel.}
\opt{D}{There exist some $\theta$ for which $S=0$ represents no real locus.}
\end{qoptions}
\qn{2} Let A,B,C be $3$ sets such that

$A=\left\{(x,y):\dfrac{x}{\cos\theta}=\dfrac{y}{\sin\theta}=2,\theta\in R\right\}$

$B=\left\{(x,y):\dfrac{x+1}{\cos2205^\circ}=\dfrac{y}{\sin2205^\circ}=r,r\in R\right\}$

$C=\left\{(x,y):(x-\sqrt3)^2+(y-1)^2\leqslant a^2;a\in R\right\}$

Let $C_1,C_2$ be the curves represented by A and B respectively while R be the region represented by C then which of following statement is/are correct
\begin{qoptions}
\opt{A}{A points $P(x_1,y_1)$ and $Q(x_2,y_2)$ be such that $P,Q\in A\cap B$, then point of intersection of tangents drawn to $C_1$ at $P$ and $Q$ is $(2,3)$.}
\opt{B}{A points $P(x_1,y_1)$ and $Q(x_2,y_2)$ be such that $P,Q\in A\cap B$, then point of intersection of tangents drawn to $C_1$ at $P$ and $Q$ is $(-4,4)$.}
\opt{C}{Minimum value of `$a$' for which $A\cap C=A$ is $\alpha$, then length of curve $C_2$ lying in R corresponding to $a=\alpha$, is $\sqrt{58}$}
\opt{D}{Minimum value of `$a$' for which $A\cap C=A$ is $\alpha$, then length of curve $C_2$ lying in R corresponding to $a=\alpha$, is $\sqrt{57}$}
\end{qoptions}
\qn{3} The sum of $99^{\mathrm{th}}$ and $100^{\mathrm{th}}$ term of the sequence $1,2,4,6,7,10,10,14,\ldots..$ is divisible by
\begin{qoptions}
\opt{A}{$2$}
\opt{B}{$53$}
\opt{C}{$109$}
\opt{D}{$173$}
\end{qoptions}
\qn{4} In a game a person P at point $(x,y)$ can jump to any one of the following points relatively: $(x+1,y)$, $(x+2,y)$, $(x,y+1)$ or $(x,y+2)$. If P is at $(0,0)$, then he can reach
\begin{qoptions}
\opt{A}{$(2,2)$ in $14$ ways}
\opt{B}{$(2,3)$ in $32$ ways}
\opt{C}{$(2,2)$ in $9$ ways}
\opt{D}{$(2,3)$ in $26$ ways}
\end{qoptions}
\qn{5} Let $x=\left(7+4\sqrt3\right)^n$, where $n$ is a natural number, then -
\begin{qoptions}
\opt{A}{$[x]$ is an even integer (where $[.]$ denotes greatest integer function)}
\opt{B}{$[x]$ is an odd integer (where $[.]$ denotes greatest integer function)}
\opt{C}{$(1-\{x\})(x)=1$ (where $\{.\}$ denotes fractional part function)}
\opt{D}{$x\{x\}=1$ (where $\{.\}$ denotes fractional part function)}
\end{qoptions}
\qn{6} The set of all values of $x$ in $[0,\pi]$ for which the function $f(x)=\sqrt{1-2\sin3x}$ takes real values is-
\begin{qoptions}
\opt{A}{$\left[0,\dfrac\pi6\right]\cup\left[\dfrac{5\pi}{6},\pi\right]$}
\opt{B}{$\left[0,\dfrac\pi9\right]\cup\left[\dfrac{2\pi}{9},\dfrac{7\pi}{9}\right]\cup\left[\dfrac{8\pi}{9},\pi\right]$}
\opt{C}{$\left[0,\dfrac{\pi}{18}\right]\cup\left[\dfrac{5\pi}{18},\dfrac{13\pi}{18}\right]\cup\left[\dfrac{17\pi}{18},\pi\right]$}
\opt{D}{$\left[0,\dfrac{13\pi}{18}\right]\cup\left[\dfrac{17\pi}{18},\pi\right]$}
\end{qoptions}
\qn{7} Sum of the solutions of the equation $[x^2]-2x+1=0$ is (where $[.]$ denotes greatest integer function)
\begin{qoptions}
\opt{A}{$\dfrac12$}
\opt{B}{$2$}
\opt{C}{$3$}
\opt{D}{$\dfrac32$}
\end{qoptions}
\qn{8} The range of the function $f(x)=\dfrac{\sqrt{1-x^2}}{1+|x|}$ is-
\begin{qoptions}
\opt{A}{$[0,1]$}
\opt{B}{$\left[0,\dfrac{1}{\sqrt2}\right]$}
\opt{C}{$\left[0,\dfrac12\right]$}
\opt{D}{$\left[0,\dfrac{\sqrt3}{2}\right]$}
\end{qoptions}
\sect{SECTION-I(ii)}
\qn{9} Answer the following by appropriately matching the lists based on the information given in the paragraph

$\alpha_1,\alpha_2,\alpha_3,\alpha_4$ and $\alpha_5$ are all roots of $x^5-4x^3-3x^2+2x-2=0$.

There are some expression given in List-I whose values are given in List-II below :

\begin{center}
{\footnotesize\setlength{\tabcolsep}{2pt}\begin{tabular}{|c|p{4.6cm}|c|p{2.4cm}|}\hline
 & \centering List-I & & \centering List-II\tabularnewline\hline
(I) & $\alpha_1^2+\alpha_2^2+\alpha_3^2+\alpha_4^2+\alpha_5^2$ & (P) & $1$\\\hline
(II) & $|(1-\alpha_1^2)(1-\alpha_2^2)\ldots.(1-\alpha_5^2)|$ & (Q) & $8$\\\hline
(III) & $|(\alpha_1^3-4\alpha_1-3)(\alpha_2^3-4\alpha_2-3)\ldots(\alpha_5^2-4\alpha_5-3)|$ & (R) & $12$\\\hline
(IV) & $\dfrac{1}{\alpha_1}+\dfrac{1}{\alpha_2}+\ldots\dfrac{1}{\alpha_5}$ & (S) & $24$\\\hline
 & & (T) & $36$\\\hline
 & & (U) & $48$\\\hline
\end{tabular}}
\end{center}
Which of the following is the only correct combination ?
\begin{qoptions}
\opt{A}{I $\to$ R}
\opt{B}{II $\to$ S}
\opt{C}{III $\to$ T}
\opt{D}{IV $\to$ Q}
\end{qoptions}
\qn{10} Answer the following by appropriately matching the lists based on the information given in the paragraph

$\alpha_1,\alpha_2,\alpha_3,\alpha_4$ and $\alpha_5$ are all roots of $x^5-4x^3-3x^2+2x-2=0$.

There are some expression given in List-I whose values are given in List-II below :

\begin{center}
{\footnotesize\setlength{\tabcolsep}{2pt}\begin{tabular}{|c|p{4.6cm}|c|p{2.4cm}|}\hline
 & \centering List-I & & \centering List-II\tabularnewline\hline
(I) & $\alpha_1^2+\alpha_2^2+\alpha_3^2+\alpha_4^2+\alpha_5^2$ & (P) & $1$\\\hline
(II) & $|(1-\alpha_1^2)(1-\alpha_2^2)\ldots.(1-\alpha_5^2)|$ & (Q) & $8$\\\hline
(III) & $|(\alpha_1^3-4\alpha_1-3)(\alpha_2^3-4\alpha_2-3)\ldots(\alpha_5^2-4\alpha_5-3)|$ & (R) & $12$\\\hline
(IV) & $\dfrac{1}{\alpha_1}+\dfrac{1}{\alpha_2}+\ldots\dfrac{1}{\alpha_5}$ & (S) & $24$\\\hline
 & & (T) & $36$\\\hline
 & & (U) & $48$\\\hline
\end{tabular}}
\end{center}
Which of the following is the only correct combination ?
\begin{qoptions}
\opt{A}{I $\to$ R}
\opt{B}{II $\to$ T}
\opt{C}{III $\to$ U}
\opt{D}{IV $\to$ R}
\end{qoptions}
\qn{11} Answer the following by appropriately matching the lists based on the information given in the paragraph

List-I contains some trigonometric equations and List-II provides number of solutions in given interval.

\begin{center}
{\footnotesize\setlength{\tabcolsep}{2pt}\begin{tabular}{|c|p{4.6cm}|c|p{2.4cm}|}\hline
 & \centering List-I & & \centering List-II\tabularnewline\hline
(I) & $2\cot2x-3\cot3x=\tan2x$ & (P) & $0$ in $\left(0,\frac\pi2\right)$\\\hline
(II) & $\sin^{10}x+\cos^{10}x=\dfrac{29}{16}\cos^42x$ & (Q) & $2$ in $(0,\pi)$\\\hline
(III) & $\cos^42x+2\sin^22x=17(\cos x+\sin x)^8$ & (R) & $4$ in $(0,\pi)$\\\hline
(IV) & $\sin\left(x+\dfrac\pi9\right)=\cos\left(2x-\dfrac{59\pi}{180}\right)$ & (S) & $4$ in $(0,2\pi)$\\\hline
 & & (T) & $2$ in $\left(-\frac\pi2,\frac\pi2\right)$\\\hline
 & & (U) & $2$ in $\left(\frac\pi2,\pi\right)$\\\hline
\end{tabular}}
\end{center}
Which of the following is the only correct combination ?
\begin{qoptions}
\opt{A}{I $\to$ T}
\opt{B}{II $\to$ R}
\opt{C}{III $\to$ P}
\opt{D}{IV $\to$ R}
\end{qoptions}
\qn{12} Answer the following by appropriately matching the lists based on the information given in the paragraph

List-I contains some trigonometric equations and List-II provides number of solutions in given interval.

\begin{center}
{\footnotesize\setlength{\tabcolsep}{2pt}\begin{tabular}{|c|p{4.6cm}|c|p{2.4cm}|}\hline
 & \centering List-I & & \centering List-II\tabularnewline\hline
(I) & $2\cot2x-3\cot3x=\tan2x$ & (P) & $0$ in $\left(0,\frac\pi2\right)$\\\hline
(II) & $\sin^{10}x+\cos^{10}x=\dfrac{29}{16}\cos^42x$ & (Q) & $2$ in $(0,\pi)$\\\hline
(III) & $\cos^42x+2\sin^22x=17(\cos x+\sin x)^8$ & (R) & $4$ in $(0,\pi)$\\\hline
(IV) & $\sin\left(x+\dfrac\pi9\right)=\cos\left(2x-\dfrac{59\pi}{180}\right)$ & (S) & $4$ in $(0,2\pi)$\\\hline
 & & (T) & $2$ in $\left(-\frac\pi2,\frac\pi2\right)$\\\hline
 & & (U) & $2$ in $\left(\frac\pi2,\pi\right)$\\\hline
\end{tabular}}
\end{center}
Which of the following is the only correct combination ?
\begin{qoptions}
\opt{A}{I $\to$ R}
\opt{B}{II $\to$ T}
\opt{C}{III $\to$ U}
\opt{D}{IV $\to$ R}
\end{qoptions}
\sect{SECTION-II}
\qn{1} If a normal chord AB to the hyperbola $xy=c^2$, normal at A, subtends an angle $\theta$ at centre and $\angle OAB=\alpha$, ($O$ is origin) then $\left|\dfrac{\tan\alpha}{\tan\theta}\right|$ is
\qn{2} Let $P$ be any point on ellipse $3x^2+4y^2=12$ and $S,S'$ are its foci then the locus of the centroid of triangle $PSS'$ is a conic $C$. Area of quadrilateral formed by tangents at the extremities of latus-rectum of the conic $C$ equals
\qn{3} If $y^2=4x$ and $y^2=(x-k)$ have a common normal other than $x$-axis, then the range of values of $k$ is $(p,\infty)$ ; where $p$ is equal to
\qn{4} From a point $P(0,5)$ on the circle $x^2+y^2=25$ two equal chords $PA$ and $PB$ of length $8$ units are drawn. The equation of $AB$ is $y=k$, where $k$ is
\qn{5} The number of real solutions of the equation $\left(\dfrac{\log_{10}x}{3}\right)^{\log_{10}^2x-3\log_{10}x+3}=\log_{10}\sqrt[3]{x}$, is
\qn{6} The number of solutions of the equation $\cos^{12}x+3\cos^4x\sin^3x=\sin^9x-1$ in $x\in[0,100\pi]$ is
\end{multicols}
\end{document}
